% =========================================================
% Proof: Finite Intersections of Open Sets are Open
% Mode: standard
% =========================================================

\subsection*{Finite Intersections of Open Sets are Open}
\label{prf:finite-intersection-open}

\begin{proof}
\Claim Let $(X,d)$ be a metric space and let
$U_1,\dots,U_n$ be open subsets of $X$.
Then $\bigcap_{k=1}^{n} U_k$ is open.

\Given A metric space $(X,d)$ and open sets $U_1,\dots,U_n \subseteq X$.

\Goal To show that $U := \bigcap_{k=1}^{n} U_k$ is open.

\Strategy Use the definition of open set:
take an arbitrary point in the intersection and use the openness
of each $U_k$ to construct a ball contained in all of them.

\medskip

Let $x \in U$ be arbitrary. Then $x \in U_k$ for each
$k = 1,\dots,n$.

Since each $U_k$ is open and $x \in U_k$, by the definition
of open set there exists $r_k > 0$ such that
\[
B(x,r_k) \subseteq U_k .
\]

Define
\[
r := \min\{r_1,r_2,\dots,r_n\}.
\]
Since the set $\{r_1,\dots,r_n\}$ is finite and each $r_k>0$,
we have $r>0$.

Let $y \in B(x,r)$. Then $d(x,y) < r \le r_k$ for each
$k = 1,\dots,n$. Hence $y \in B(x,r_k) \subseteq U_k$
for every $k$.

Therefore $y \in \bigcap_{k=1}^{n} U_k = U$.
Thus
\[
B(x,r) \subseteq U .
\]

Since $x$ was arbitrary, every point of $U$ has an open ball
contained in $U$. Therefore $U$ is open.
\end{proof}

\begin{remark*}[Dependencies]
\begin{itemize}
  \item Definition of open ball.
  \item Definition of open set.
  \item (Key idea) Use a minimum of finitely many radii.
\end{itemize}
\end{remark*}




% =========================================================
% TikZ: Finite Intersection of Open Sets is Open (U1, U2)
% Matches the picture: two big overlapping sets + three concentric balls
% "50% bigger" achieved via scale=1.5
% =========================================================

\begin{tikzpicture}[scale=1.5, line cap=round, line join=round]

  %--- Colors (close to your screenshot)
  \definecolor{Uone}{RGB}{58,103,190}   % blue stroke
  \definecolor{Utwo}{RGB}{176,94,52}    % brown/orange stroke

  %--- Centers/radii for the two sets (tweak overlap by changing their centers)
  \coordinate (c1) at (-2.3,0.0);   % move right/left to adjust overlap
\coordinate (c2) at ( 1.9,0.0);
   % move left/right to adjust overlap
  \def\Rset{2.25}

  %--- Point x in the overlap
  \coordinate (x) at (0.0,0.0);

  %--- Radii for the balls around x
  \def\rmin{1.05}   % r = min{r1,r2}
  \def\rone{1.25}   % r1
  \def\rtwo{1.70}   % r2

  %--- Title
  \node at (0,2.85) {$r=\min\{r_1,r_2\}$};

  %--- Fill the open sets (light)
  \fill[Uone!12] (c1) circle (\Rset);
  \fill[Utwo!12] (c2) circle (\Rset);

  %--- Boundaries of U1, U2
  \draw[Uone, line width=0.9pt] (c1) circle (\Rset);
  \draw[Utwo, line width=0.9pt] (c2) circle (\Rset);

  %--- Labels for U1, U2
  \node[Uone] at ($(c1)+(0,1.35)$) {$U_1$};
  \node[Utwo] at ($(c2)+(0,1.35)$) {$U_2$};

  %--- Balls around x
  \draw[black, dashed, line width=0.7pt] (x) circle (\rtwo);
  \draw[black, dashed, line width=0.7pt] (x) circle (\rone);
  \draw[black, line width=0.8pt]         (x) circle (\rmin);

  %--- Point x and a radius segment labeled r
  \fill[black] (x) circle (1.2pt);
  \node[above right] at (x) {$x$};

  \draw[black, line width=0.5pt] (x) -- ++(25:\rmin);
  \node at ($(x)+(25:0.62*\rmin)$) {\scriptsize $r$};

  %--- Labels for the balls (positions chosen to mimic your screenshot)
  \node at (-0.95, 0.20) {$B(x,r_1)$};
  \node at ( 0.95,-0.15) {$B(x,r_2)$};
  \node at ( 0.00,-1.15) {$B(x,r)$};

  %--- Intersection label
  \node at (0.00,-0.65) {$U_1\cap U_2$};

\end{tikzpicture}













